% ============================================================================
%  SECCIÓN 1.13: POBLACIÓN Y MUESTRA
% ============================================================================

\section{Poblaci\'on y muestra}

La poblaci\'on es el conjunto de todos los individuos, objetos y eventos
sujetos a investigaci\'on, de los cuales se desea estudiar una propiedad o
caracter\'istica observable; la muestra es un subconjunto o una parte de la
poblaci\'on en la que se lleva a cabo la investigaci\'on, con el fin de
generalizar los hallazgos al todo \citep{hernandez2014}.

Para el presente proyecto, la poblaci\'on de estudio est\'a constituida por dos
grupos: el \textbf{p\'ublico general} potencial usuario del sistema
(asistentes a eventos deportivos, aficionados y deportistas) y los
\textbf{gestores de los escenarios deportivos} del pa\'is.

\subsection{Poblaci\'on y muestra de los gestores de escenarios}

En el caso de los gestores de escenarios deportivos, se considera necesario
estudiar una muestra de los escenarios de las ciudades de La Paz,
Cochabamba y Santa Cruz de la Sierra, donde se concentra la mayor parte de la
oferta deportiva nacional. Se utilizar\'a un muestreo no probabil\'istico por
conveniencia, seleccionando los escenarios de mayor actividad deportiva en
cada ciudad, debido a la accesibilidad y a la disposici\'on de las
instituciones a participar en el estudio.

\subsection{Poblaci\'on y muestra del p\'ublico general}

En el caso del p\'ublico general, la muestra se determinar\'a mediante la
f\'ormula para poblaciones finitas, considerando la poblaci\'on de las tres
ciudades del piloto, con un nivel de confianza del 95\,\% y un error m\'aximo
del 5\,\%. El muestreo ser\'a probabil\'istico aleatorio simple, con aplicaci\'on
del cuestionario mediante medios digitales.

\subsection{Justificaci\'on de la muestra}

Se justifica el estudio de una muestra porque la aplicaci\'on del instrumento a
toda la poblaci\'on resultar\'ia costosa e innecesaria, y las caracter\'isticas de
la poblaci\'on son homog\'eneas respecto a las variables de inter\'es del estudio.
La muestra seleccionada cumple los requisitos de ser confiable y
representativa, permitiendo extrapolar los resultados obtenidos a la
poblaci\'on en su conjunto.
